
\documentclass[12pt,a4paper]{article}

\usepackage[utf8]{inputenc}
\usepackage[T1]{fontenc}
\usepackage[spanish,es-tabla]{babel}
\usepackage{lmodern}
\usepackage[a4paper,margin=2.2cm]{geometry}
\usepackage{array}
\usepackage{longtable}
\usepackage{microtype}
\usepackage{hyperref}

\hypersetup{
    colorlinks=true,
    linkcolor=blue,
    urlcolor=blue
}

\setlength{\parindent}{0pt}
\setlength{\parskip}{6pt}

\begin{document}

\begin{titlepage}
\centering

{\bfseries UNIVERSIDAD TÉCNICA ESTATAL DE QUEVEDO\par}
\vspace{0.7cm}

FACULTAD DE CIENCIAS DE LA COMPUTACIÓN\par
CARRERA DE SOFTWARE\par
INGENIERÍA DE REQUERIMIENTOS\par

\vfill

{\Large\bfseries GA-PRAC-01 - FORM\par}
\vspace{0.5cm}

{\large Comparación entre modelos de proceso de requisitos en escenarios reales\par}

\vspace{1cm}

\textbf{Integrantes:}\par
Córdova Carreño Mayra Lucila\par
Herrera Morán Gilmar Javier\par
Rivas Rodero Jordy Fabian\par

\vfill

Noviembre de 2026\par

\end{titlepage}

\tableofcontents
\clearpage

% C1: Caracterización de modelos
% Responsable: Córdova Carreño Mayra Lucila
\input{01-Cordova/c1-caracterizacion.tex}

% C2: Matrices comparativas
% Responsable: Córdova Carreño Mayra Lucila
\input{01-Cordova/c2-escenario-01.tex}
\input{01-Cordova/c2-escenario-02.tex}
\input{01-Cordova/c2-escenario-03.tex}
\input{01-Cordova/c2-escenario-04.tex}

% C3: Actores del proceso y framework de IR
% Responsable: Herrera Morán Gilmar Javier
\input{02-Herrera/c3-actores-framework.tex}

% C4 y C5:
% Pendientes de integrar cuando Rivas suba su archivo.

% Bibliografía consolidada del documento (C3 y C4-C5).
% Si Rivas integra su archivo, debe quitar su propio
% bloque thebibliography para no duplicar las referencias.
\begin{thebibliography}{9}
\bibitem{Pohl2025} Pohl, K. (2025). \emph{Requirements Engineering: Fundamentals, Principles, and Techniques} (2.ª ed.). Springer. ISBN 978-3-662-69204-2.
\bibitem{IREB2026} Bühne, S., Glinz, M., van Loenhoud, H. y Staal, S. (2026). \emph{CPRE Foundation Level Syllabus, versión 3.3.0}. International Requirements Engineering Board (IREB). \url{https://www.ireb.org/en/downloads/}
\bibitem{SWEBOK_v4_2024} Washizaki, H. (Ed.). (2024). \emph{Guide to the Software Engineering Body of Knowledge (SWEBOK Guide), Version 4.0}. IEEE Computer Society. \url{https://www.computer.org/education/bodies-of-knowledge/software-engineering}
\bibitem{ISO29148} ISO/IEC/IEEE. (2018). \emph{ISO/IEC/IEEE 29148:2018, Systems and software engineering — Life cycle processes — Requirements engineering} (2.ª ed.). \url{https://webstore.iec.ch/en/publication/64315}
\bibitem{Inayat2015} Inayat, I., Salim, S. S., Marczak, S., Daneva, M. y Shamshirband, S. (2015). A systematic literature review on agile requirements engineering practices and challenges. \emph{Computers in Human Behavior}, 51(B), 915--929. \url{https://doi.org/10.1016/j.chb.2014.10.046}
\bibitem{Sommerville2016} Sommerville, I. (2016). \emph{Software Engineering} (10.ª ed.). Pearson. ISBN 978-0-13-394303-0.
\end{thebibliography}

\end{document}
